\documentclass[11pt]{article}
\usepackage[a4paper,margin=22mm]{geometry}
\usepackage{fontspec}
\setmainfont{DejaVu Serif}
\setsansfont{DejaVu Sans}
\usepackage{microtype}
\usepackage{xcolor}
\usepackage{hyperref}
\usepackage{enumitem}
\usepackage{parskip}
\definecolor{Accent}{HTML}{971D32}
\definecolor{Muted}{HTML}{666666}
\hypersetup{colorlinks=true,urlcolor=Accent,linkcolor=Accent}
\setlist{leftmargin=1.6em,itemsep=0.3em,topsep=0.35em}
\setcounter{tocdepth}{2}
\renewcommand{\contentsname}{Contents}
\newcommand{\sectionrule}{\par\vspace{-0.5em}\noindent\textcolor{Accent}{\rule{\linewidth}{0.6pt}}\par}
\begin{document}

\begin{center}
{\sffamily\bfseries\color{Accent} TOEFL WORD OF THE DAY\par}
\vspace{8mm}
{\fontsize{42}{48}\selectfont\bfseries inventory\par}
\vspace{2mm}
{\Large\sffamily\color{Accent} /ˈɪnvənˌtɔːri/\par}
\vspace{2mm}
{\small\sffamily Local audio phrase: inventory\par}
{\small\sffamily Audio file: audios/inventory.mp3\par}
\vspace{2mm}
{\large\sffamily\color{Muted} countable or uncountable noun; transitive verb\par}
\vspace{5mm}
{\sffamily goods kept for use or sale \textbullet\ a detailed list of items \textbullet\ make a detailed list\par}
\vspace{6mm}
{\large Inventory can refer to the goods an organization has, a detailed list of those items, or the act of making such a list.\par}
\vspace{6mm}
\begin{minipage}{0.84\textwidth}
\itshape “Retailers often increase their inventory before major shopping seasons.”
\end{minipage}\par
\vspace{10mm}
\textbf{Date:} August 31st 2026\quad\textbullet\quad \textbf{Level:} B1--mid-B2 TOEFL
\vfill
Created and compiled by \textbf{Amir Kooshky}\par
\vspace{2mm}
\href{https://t.me/KooshkyTOEFL}{Telegram Channel}\quad\textbullet\quad
\href{https://www.instagram.com/kooshkytoefl}{Instagram}\quad\textbullet\quad
\href{https://t.me/KooshkyTOEFL\_pv}{Personal Telegram}
\end{center}
\newpage
\tableofcontents
\newpage

\section{Meanings and usage}
\sectionrule
\subsection{Meaning 1: Goods kept for use or sale}
\textbf{Part of speech:} countable or uncountable noun

\textbf{IPA:} /ˈɪnvənˌtɔːri/

\textbf{Pronunciation phrase:} inventory

\textbf{Local audio file:} audios/inventory.mp3

\textbf{Register:} neutral to formal; very common in business, economics, manufacturing, retail, and supply-chain contexts

\textbf{Definition:} the goods, materials, or products that a business or organization has available or stored for future use or sale

\textbf{In simpler words:} the products or materials a business has in stock

\textbf{Typical pattern:} inventory of + products or materials; keep or hold + inventory; in inventory

\subsubsection{Natural examples}
\begin{enumerate}
  \item The company reduced its inventory after demand for the product declined.
  \item Retailers often increase their inventory before major shopping seasons.
  \item The factory keeps enough raw materials in inventory to avoid interruptions in production.
  \item Poor inventory management can cause both shortages and unnecessary storage costs.
  \item The store has a large inventory of electronic devices.
  \item Manufacturers try to avoid keeping excessive inventory when demand is uncertain.
  \item The company sold much of its remaining inventory at a discount.
  \item Supply chain problems made it difficult for businesses to maintain adequate inventory.
\end{enumerate}

\subsubsection{Nuance and implication}
\textbf{Central nuance:} Inventory refers to goods or materials that are currently available rather than goods that have already been sold or consumed.

\textbf{Usual implication:} Businesses often need to balance having enough inventory to meet demand against the cost of storing too much.

\textbf{Compared with sales:} Inventory consists of goods a business still has available, while sales refer to goods that have already been sold.

\subsubsection{Grammar notes}
\textbf{How to use it:} Inventory can be uncountable when referring generally to a company's stock of goods, but it can also be countable when referring to a particular stock or range of items.

\textbf{What can follow it:} It often appears with verbs such as keep, hold, maintain, reduce, increase, manage, sell, or replenish.

\textbf{Common preposition:} Use in inventory or in stock for goods currently available. Use inventory of when naming the type of goods.

\textbf{Noun use:} In business English, inventory is often uncountable: The company has too much inventory. A large inventory is also possible when referring to a particular quantity or collection of stock.

\subsubsection{Useful patterns}
\textbf{Pattern:} keep or hold inventory

\textbf{Use:} Use this for goods or materials a company stores for future use or sale.

\textbf{Example:} The retailer keeps relatively little inventory in its smaller stores.

\textbf{Pattern:} inventory of + products or materials

\textbf{Use:} Use this when naming what kinds of items a business has available.

\textbf{Example:} The warehouse contains a large inventory of replacement parts.

\textbf{Pattern:} in inventory

\textbf{Use:} Use this for goods currently being held by a business or organization.

\textbf{Example:} The company still has several thousand units in inventory.

\textbf{Pattern:} reduce or increase inventory

\textbf{Use:} Use this when a business changes the quantity of goods it stores.

\textbf{Example:} The manufacturer reduced inventory after orders began to fall.

\textbf{Pattern:} inventory management

\textbf{Use:} Use this for the process of controlling how much stock a business keeps and when it replaces it.

\textbf{Example:} Effective inventory management can reduce storage costs.

\subsubsection{Common wrong usage}
\paragraph{Correction 1}
\textbf{Wrong:} The company has many inventory in its warehouse.

\textbf{Correct:} The company has a lot of inventory in its warehouse.

\textbf{Why:} When inventory means stock in general, it is often uncountable, so use a lot of rather than many.

\paragraph{Correction 2}
\textbf{Wrong:} The store increased its inventories before the holiday season.

\textbf{Correct:} The store increased its inventory before the holiday season.

\textbf{Why:} When referring generally to the store's total stock, the singular or uncountable form inventory is more natural.

\paragraph{Correction 3}
\textbf{Wrong:} The products are on inventory.

\textbf{Correct:} The products are in inventory.

\textbf{Why:} Use in inventory, not on inventory.

\subsection{Meaning 2: Detailed list of items}
\textbf{Part of speech:} countable noun

\textbf{IPA:} /ˈɪnvənˌtɔːri/

\textbf{Pronunciation phrase:} inventory

\textbf{Local audio file:} audios/inventory.mp3

\textbf{Register:} neutral to formal; common in museums, science, administration, environmental studies, archaeology, and resource management

\textbf{Definition:} a complete and detailed list of the items, resources, or materials that a person or organization has

\textbf{In simpler words:} a complete list of what is available

\textbf{Typical pattern:} an inventory of + objects, species, resources, equipment, or materials

\subsubsection{Natural examples}
\begin{enumerate}
  \item The museum created an inventory of all the artifacts in its collection.
  \item Researchers prepared an inventory of plant species found in the protected area.
  \item The laboratory maintains a detailed inventory of its chemicals and equipment.
  \item Before the building was renovated, staff compiled an inventory of its historical features.
  \item The agency developed an inventory of available water resources.
  \item Archaeologists made an inventory of the objects recovered from the site.
  \item The school keeps an inventory of computers and other electronic equipment.
  \item Officials used the inventory to determine which supplies were missing.
\end{enumerate}

\subsubsection{Nuance and implication}
\textbf{Central nuance:} An inventory is usually systematic and comprehensive rather than a quick or informal list.

\textbf{Usual implication:} The purpose is often to record what exists, where it is, how much there is, or what condition it is in.

\textbf{Compared with sample:} A sample contains only some items selected from a larger group, while an inventory aims to record all or nearly all relevant items.

\subsubsection{Grammar notes}
\textbf{How to use it:} This meaning is countable. Say an inventory for one detailed list and inventories for more than one.

\textbf{What can follow it:} It is commonly followed by of and the category being recorded.

\textbf{Common preposition:} Use an inventory of something: an inventory of equipment, an inventory of species.

\textbf{Noun use:} The singular normally needs a determiner such as an, the, or this.

\subsubsection{Useful patterns}
\textbf{Pattern:} an inventory of + items

\textbf{Use:} Use this for a detailed record of a collection or group of things.

\textbf{Example:} The researchers compiled an inventory of native plant species.

\textbf{Pattern:} create or compile an inventory

\textbf{Use:} Use this when people systematically produce a detailed list.

\textbf{Example:} The museum compiled an inventory of objects damaged in the flood.

\textbf{Pattern:} maintain an inventory

\textbf{Use:} Use this when an organization keeps a list accurate and up to date.

\textbf{Example:} The laboratory maintains an inventory of hazardous chemicals.

\textbf{Pattern:} detailed or comprehensive inventory

\textbf{Use:} Use these adjectives when emphasizing completeness or detail.

\textbf{Example:} The survey produced a comprehensive inventory of the region's wetlands.

\subsubsection{Common wrong usage}
\paragraph{Correction 1}
\textbf{Wrong:} Researchers made an inventory for the species in the forest.

\textbf{Correct:} Researchers made an inventory of the species in the forest.

\textbf{Why:} Use an inventory of something when naming the items included in the list.

\paragraph{Correction 2}
\textbf{Wrong:} The museum prepared inventory of its paintings.

\textbf{Correct:} The museum prepared an inventory of its paintings.

\textbf{Why:} This meaning is countable, so the singular needs a determiner such as an.

\paragraph{Correction 3}
\textbf{Wrong:} An inventory is just one example taken from a larger collection.

\textbf{Correct:} An inventory is a detailed list of the items in a collection.

\textbf{Why:} A sample contains only selected examples; an inventory aims to record the collection systematically.

\subsection{Meaning 3: Make a detailed list}
\textbf{Part of speech:} transitive verb

\textbf{IPA:} /ˈɪnvənˌtɔːri/

\textbf{Pronunciation phrase:} inventory

\textbf{Local audio file:} audios/inventory.mp3

\textbf{Register:} formal; common in administration, science, museums, archaeology, business, and resource management

\textbf{Definition:} to make a detailed list of the items or resources in a place or collection

\textbf{In simpler words:} make a complete list of

\textbf{Typical pattern:} inventory + objects, equipment, resources, species, or supplies

\subsubsection{Natural examples}
\begin{enumerate}
  \item Museum staff inventoried thousands of objects before moving the collection.
  \item Researchers inventoried the plant species found in each study area.
  \item The organization inventoried its equipment to determine what needed to be replaced.
  \item Archaeologists carefully inventoried every object recovered from the site.
  \item Employees spent several days inventorying the contents of the warehouse.
  \item The agency inventoried available water resources throughout the region.
  \item Before closing the laboratory, staff inventoried all remaining chemicals.
\end{enumerate}

\subsubsection{Nuance and implication}
\textbf{Central nuance:} As a verb, inventory means systematically identify and record what is present.

\textbf{Usual implication:} The process often includes counting, describing, locating, or assessing the condition of the items.

\subsubsection{Grammar notes}
\textbf{How to use it:} This verb normally needs a direct object: inventory the equipment, inventory the collection.

\textbf{What can follow it:} Common objects include equipment, supplies, objects, resources, species, assets, and collections.

\textbf{Position in a sentence:} The -ing form is inventorying, with the y kept: inventorying the supplies.

\subsubsection{Useful patterns}
\textbf{Pattern:} inventory + equipment or supplies

\textbf{Use:} Use this when systematically recording what equipment or materials are available.

\textbf{Example:} Staff inventoried the emergency supplies before the hurricane season.

\textbf{Pattern:} inventory + collection

\textbf{Use:} Use this when recording the contents of a museum, archive, or other collection.

\textbf{Example:} The museum began inventorying its entire archaeological collection.

\textbf{Pattern:} inventory + species or resources

\textbf{Use:} Use this in scientific or environmental contexts when recording what exists in an area.

\textbf{Example:} Scientists inventoried the species living in the wetland.

\subsubsection{Common wrong usage}
\paragraph{Correction 1}
\textbf{Wrong:} The researchers inventoried about the species in the forest.

\textbf{Correct:} The researchers inventoried the species in the forest.

\textbf{Why:} Inventory is transitive as a verb, so put the thing being recorded directly after it.

\paragraph{Correction 2}
\textbf{Wrong:} The staff was inventory the equipment.

\textbf{Correct:} The staff was inventorying the equipment.

\textbf{Why:} After was, use the -ing form inventorying.

\section{Word family}
\sectionrule
\subsection{inventorial}
\textbf{Part of speech:} adjective

\textbf{IPA:} /ˌɪnvənˈtɔːriəl/

\textbf{Definition:} relating to an inventory or the process of making one

\textbf{Relationship to the headword:} This adjective describes records or information connected with an inventory.

\textbf{Grammar pattern:} inventorial + record, data, or description

\textbf{Usage note:} This word is uncommon and mainly appears in specialized formal contexts.

\subsubsection{Examples}
\begin{enumerate}
  \item The archive contains inventorial records describing the collection.
  \item Researchers used inventorial data to compare the available resources.
\end{enumerate}

\section{Common collocations}
\sectionrule
\subsection{inventory management}
\textbf{Applies to:} Goods kept for use or sale

\textbf{Meaning:} the process of controlling how much stock a business keeps

\textbf{Example:} Better inventory management helped the company reduce storage costs.

\subsection{inventory levels}
\textbf{Applies to:} Goods kept for use or sale

\textbf{Meaning:} the amount of stock currently available

\textbf{Example:} Retailers increased inventory levels before the holiday season.

\subsection{reduce inventory}
\textbf{Applies to:} Goods kept for use or sale

\textbf{Meaning:} decrease the amount of goods or materials being stored

\textbf{Example:} The manufacturer reduced inventory when demand began to decline.

\subsection{increase inventory}
\textbf{Applies to:} Goods kept for use or sale

\textbf{Meaning:} increase the amount of goods being held

\textbf{Example:} The company increased inventory to prepare for stronger demand.

\subsection{excess inventory}
\textbf{Applies to:} Goods kept for use or sale

\textbf{Meaning:} more stock than is currently needed

\textbf{Example:} The retailer offered discounts to sell excess inventory.

\subsection{inventory of goods}
\textbf{Applies to:} Goods kept for use or sale

\textbf{Meaning:} a stock of products available for sale or use

\textbf{Example:} The warehouse contains a large inventory of imported goods.

\subsection{in inventory}
\textbf{Applies to:} Goods kept for use or sale

\textbf{Meaning:} currently held as stock

\textbf{Example:} The company still has thousands of unsold units in inventory.

\subsection{inventory of species}
\textbf{Applies to:} Detailed list of items

\textbf{Meaning:} a systematic list of species present in an area

\textbf{Example:} The survey produced an inventory of species found in the wetland.

\subsection{inventory of resources}
\textbf{Applies to:} Detailed list of items

\textbf{Meaning:} a detailed record of available resources

\textbf{Example:} The agency prepared an inventory of available water resources.

\subsection{compile an inventory}
\textbf{Applies to:} Detailed list of items

\textbf{Meaning:} systematically create a detailed list

\textbf{Example:} Researchers compiled an inventory of historical buildings in the city.

\subsection{maintain an inventory}
\textbf{Applies to:} Detailed list of items

\textbf{Meaning:} keep a detailed list accurate and up to date

\textbf{Example:} The laboratory maintains an inventory of all hazardous substances.

\section{Synonyms and antonyms}
\sectionrule
\subsection{Close synonyms}
\subsubsection{stock}
\textbf{Part of speech:} countable or uncountable noun

\textbf{Applies to:} Goods kept for use or sale

\textbf{Shared meaning:} Both can refer to goods that a business currently has available.

\textbf{Difference:} Stock is the common everyday business word for goods available for sale or use. Inventory is especially common in American business English and often emphasizes the quantity being stored and managed.

\textbf{Example:} The store does not have enough stock to meet current demand.

\subsubsection{supply}
\textbf{Part of speech:} countable noun

\textbf{Applies to:} Goods kept for use or sale

\textbf{Shared meaning:} Both can refer to materials kept for future use.

\textbf{Difference:} A supply is an amount of something available for use, especially something that may be consumed. Inventory more often refers to goods or materials systematically held and tracked by an organization.

\textbf{Example:} The hospital maintains an emergency supply of medicine.

\subsubsection{catalog}
\textbf{Part of speech:} countable noun

\textbf{Applies to:} Detailed list of items

\textbf{Shared meaning:} Both can be organized lists of items.

\textbf{Difference:} A catalog is an organized list, often with descriptions and sometimes intended for users or customers. An inventory is mainly a record of what items actually exist or are currently held.

\textbf{Example:} The library published an online catalog of its rare books.

\subsubsection{list}
\textbf{Part of speech:} countable noun

\textbf{Applies to:} Detailed list of items

\textbf{Shared meaning:} Both can record a group of items.

\textbf{Difference:} List is the general word. Inventory usually suggests a systematic and relatively complete record of possessions, resources, or items in a collection.

\textbf{Example:} The researcher made a list of the samples collected that day.

\subsection{Direct or useful antonyms}
\subsubsection{shortage}
\textbf{Part of speech:} countable noun

\textbf{Applies to:} Goods kept for use or sale

\textbf{Opposition:} Inventory is the stock of goods or materials available, while a shortage is a situation in which there is not enough of something to meet demand.

\textbf{Usage note:} This is a useful contrast rather than a perfect dictionary antonym.

\textbf{Example:} The factory experienced a shortage of essential components.

\section{Words learners may confuse}
\sectionrule
\subsection{stock}
\textbf{Part of speech:} noun

\textbf{Why learners confuse them:} The two words overlap strongly when referring to goods available for sale.

\textbf{Key difference:} Inventory can mean the goods a business holds, especially in American business English. Stock can have the same meaning, but it also has several unrelated meanings, including shares in a company.

\textbf{inventory:} The retailer has too much inventory in its warehouses.

\textbf{stock:} The retailer has too much stock in its warehouses.

\subsection{catalog}
\textbf{Part of speech:} countable noun

\textbf{Why learners confuse them:} Both can refer to organized lists of objects.

\textbf{Key difference:} An inventory records what items are actually present or owned. A catalog is an organized list or presentation of items and may include things that are available, described, or offered rather than physically present.

\textbf{inventory:} The museum completed an inventory of all objects in storage.

\textbf{catalog:} The museum published a catalog describing the exhibition.

\subsection{invent}
\textbf{Part of speech:} verb

\textbf{Why learners confuse them:} The words begin similarly but have different modern meanings and uses.

\textbf{Key difference:} Inventory refers to stock or a detailed list, and as a verb means make such a list. Invent means create or design something new.

\textbf{inventory:} Staff inventoried all the equipment in the laboratory.

\textbf{invent:} The engineer invented a new type of sensor.

\section{Etymology}
\sectionrule
\textbf{Origin:} Inventory comes through Medieval Latin and French from a word referring to a detailed list of goods or possessions.

\textbf{Development:} The meaning expanded from the written list itself to the goods or stock recorded on that list.

\textbf{Memory link:} An inventory originally focuses on recording what you have; today the word can mean either the list or the actual goods on that list.

\section{Practice exercises}
\sectionrule
\subsection{Correct or incorrect?}
Decide whether each sentence is correct. Correct the inaccurate ones.

\begin{enumerate}
  \item The retailer increased its inventory before the holiday season.
  \item The company has many inventory in its warehouse.
  \item Researchers compiled an inventory of the species found in the wetland.
  \item The museum made an inventory for all its artifacts.
  \item Staff inventoried the laboratory equipment before the move.
\end{enumerate}

\subsection{Collocation practice}
Complete each sentence with a natural collocation.

\begin{enumerate}
  \item Effective inventory \_\_\_\_\_ can help a company avoid both shortages and excessive storage costs.
  \item The museum compiled an inventory \_\_\_\_\_ all the objects in its collection.
  \item The retailer reduced \_\_\_\_\_ after sales declined.
  \item Scientists prepared a detailed inventory of plant \_\_\_\_\_ in the protected area.
\end{enumerate}

\subsection{Complete answer key}
\subsubsection{Correct or incorrect?}
\begin{enumerate}
  \item \textbf{Correct} \emph{Inventory naturally refers to goods a business keeps available for sale.}
  \item \textbf{Incorrect}. The company has a lot of inventory in its warehouse. \emph{When inventory means stock in general, it is commonly uncountable.}
  \item \textbf{Correct} \emph{An inventory of something is a systematic list of the items present.}
  \item \textbf{Incorrect}. The museum made an inventory of all its artifacts. \emph{Use an inventory of when naming the items included in the list.}
  \item \textbf{Correct} \emph{As a transitive verb, inventory means make a detailed record of the items.}
\end{enumerate}

\subsubsection{Collocation practice}
\begin{enumerate}
  \item \textbf{management.} Effective inventory management can help a company avoid both shortages and excessive storage costs. \emph{Inventory management is the process of controlling and organizing a company's stock.}
  \item \textbf{of.} The museum compiled an inventory of all the objects in its collection. \emph{Use an inventory of something when naming what the list contains.}
  \item \textbf{inventory.} The retailer reduced inventory after sales declined. \emph{Reduce inventory means decrease the amount of stock being held.}
  \item \textbf{species.} Scientists prepared a detailed inventory of plant species in the protected area. \emph{An inventory of species is a systematic record of the species present in an area.}
\end{enumerate}

\vfill
\begin{center}
\small\color{Muted} Created and compiled by \textbf{Amir Kooshky}\\
\href{https://t.me/KooshkyTOEFL}{Telegram Channel}\quad\textbullet\quad
\href{https://www.instagram.com/kooshkytoefl}{Instagram}\quad\textbullet\quad
\href{https://t.me/KooshkyTOEFL\_pv}{Personal Telegram}
\end{center}
\end{document}
